% =====================================================================
%  CHAPITRE 10 — GRAMMAIRE B1 : MODAUX, CONDITIONNELS, PASSIF, RAPPORTÉ
% =====================================================================
\chapter{Structures intermédiaires : modaux, conditions, passif, récit}

Ce chapitre contient les structures qui font passer un discours de « scolaire » à « adulte ».
Chacune est courte à comprendre et longue à automatiser : apprenez-les en blocs, pas en règles.

\section{Unité 1 --- Les verbes modaux}

\begin{regle}[Ce qu'est un modal]
Un modal (\en{can, could, may, might, must, shall, should, will, would}) ne change
\textbf{jamais} : pas de \en{-s}, pas de \en{-ing}, pas de \en{to} après lui, et il forme
seul la négation et la question.\\
\en{She can swim.} \quad \en{She can't swim.} \quad \en{Can she swim?}\\
\textbf{Structure :} sujet + modal + \textbf{base verbale}.
\end{regle}

\begin{tableau}[Le grand tableau des modaux et de leurs sens]
\footnotesize
\begin{longtable}{@{}p{2.6cm}p{4.0cm}p{4.6cm}p{2.6cm}@{}}
\entete{Modal} & \entete{Sens} & \entete{Exemple} & \entete{Traduction} \tabularnewline
\endfirsthead
\entete{Modal} & \entete{Sens} & \entete{Exemple} & \entete{Traduction} \tabularnewline
\endhead
\hline
\en{can} & capacité, possibilité, permission informelle & \en{I can drive. Can I help?} & je peux, je sais \tabularnewline
\en{could} & capacité passée, demande polie, possibilité faible & \en{I could swim at 5. Could you help me?} & je pouvais, pourriez-vous \tabularnewline
\en{be able to} & capacité (remplace \en{can} aux temps manquants) & \en{I will be able to come. I haven't been able to sleep.} & je pourrai, je n'ai pas pu \tabularnewline
\en{may} & permission formelle, possibilité moyenne & \en{May I come in? It may rain.} & puis-je, il se peut \tabularnewline
\en{might} & possibilité faible (20-30\,\%) & \en{She might be late.} & elle sera peut-être en retard \tabularnewline
\en{must} & obligation forte venant de moi, déduction logique & \en{I must study. He must be tired.} & je dois, il doit être \tabularnewline
\en{have to} & obligation venant de l'extérieur (règle, loi) & \en{I have to wear a uniform.} & je dois (obligé par les circonstances) \tabularnewline
\en{must not} & \textbf{interdiction} & \en{You mustn't smoke here.} & tu ne dois surtout pas \tabularnewline
\en{don't have to} & \textbf{absence d'obligation} & \en{You don't have to pay.} & tu n'es pas obligé (mais tu peux) \tabularnewline
\en{should / ought to} & conseil, recommandation & \en{You should rest. You ought to call her.} & tu devrais \tabularnewline
\en{will} & futur, décision, promesse, habitude & \en{I'll do it. He will always help.} & je ferai \tabularnewline
\en{would} & conditionnel, habitude passée, politesse & \en{I would go if I could. Would you like tea?} & j'irais, voudriez-vous \tabularnewline
\en{shall} & offre ou suggestion (surtout \en{I/we}, britannique) & \en{Shall we dance? Shall I call him?} & veux-tu que\ldots ? \tabularnewline
\en{need} & nécessité (surtout négatif et interrogatif) & \en{You needn't worry. Need I come?} & tu n'as pas besoin \tabularnewline
\en{dare} & oser (littéraire) & \en{How dare you! I daren't tell her.} & comment oses-tu \tabularnewline
\end{longtable}
\end{tableau}

\begin{piege}[\en{mustn't} ne veut pas dire « ne pas être obligé »]
\en{You \textbf{mustn't} tell anyone.} = il est \textbf{interdit} de le dire.\\
\en{You \textbf{don't have to} tell anyone.} = tu n'es \textbf{pas obligé} de le dire.\\
Le français « tu ne dois pas » est ambigu ; l'anglais non. C'est une faute de sens très
coûteuse (consignes de sécurité, règles de travail, examens).
\end{piege}

\begin{piege}[\en{*I can to go}, \en{*he cans}, \en{*I am can}]
Trois fautes de débutant qui persistent : pas de \en{to} après un modal, pas de \en{-s},
pas de \en{be}. Et pour les temps où \en{can} n'existe pas, utilisez \en{be able to} :
\en{I will be able to} (futur), \en{I have been able to} (\en{perfect}),
\en{I would like to be able to} (conditionnel).
\end{piege}

\subsection{Les modaux de déduction : « ça doit être », « ça ne peut pas être »}
\begin{tableau}[Échelle de certitude]
\begin{tabularx}{\linewidth}{@{}p{3.2cm}p{2.4cm}X@{}}
\entete{Structure} & \entete{Certitude} & \entete{Exemple} \tabularnewline
\hline
\en{must} & 95\,\% (presque sûr) & \en{He must be at home --- his car is here.} \textit{(Il doit être chez lui.)} \tabularnewline
\en{will} & 90\,\% & \en{That will be the postman.} \textit{(Ce sera le facteur.)} \tabularnewline
\en{should} & 75\,\% & \en{She should be there by now.} \textit{(Elle devrait y être.)} \tabularnewline
\en{may} & 50\,\% & \en{It may be true.} \textit{(C'est peut-être vrai.)} \tabularnewline
\en{might / could} & 30\,\% & \en{He might be sleeping.} \textit{(Il dort peut-être.)} \tabularnewline
\en{can't} & 5\,\% (impossible) & \en{That can't be Ali --- he's in Paris!} \textit{(Ça ne peut pas être Ali.)} \tabularnewline
\end{tabularx}
\end{tableau}

\subsection{Les modaux au passé : \en{modal} + \en{have} + participe passé}
\begin{tableau}[Le regret, la déduction passée, la possibilité manquée]
\begin{tabularx}{\linewidth}{@{}p{3.6cm}p{4.4cm}X@{}}
\entete{Structure} & \entete{Sens} & \entete{Exemple} \tabularnewline
\hline
\en{must have done} & déduction sur le passé (c'est forcément arrivé) & \en{She must have forgotten.} \textit{(Elle a dû oublier.)} \tabularnewline
\en{can't have done} & impossibilité passée & \en{He can't have left already.} \textit{(Il ne peut pas être déjà parti.)} \tabularnewline
\en{may/might have done} & possibilité passée & \en{They might have missed the bus.} \textit{(Ils ont peut-être raté le bus.)} \tabularnewline
\en{should have done} & regret, reproche (j'aurais dû) & \en{I should have studied more.} \textit{(J'aurais dû étudier davantage.)} \tabularnewline
\en{shouldn't have done} & reproche inverse (je n'aurais pas dû) & \en{You shouldn't have said that.} \textit{(Tu n'aurais pas dû dire ça.)} \tabularnewline
\en{could have done} & possibilité non réalisée & \en{I could have won, but I fell.} \textit{(J'aurais pu gagner, mais je suis tombé.)} \tabularnewline
\en{needn't have done} & action faite alors qu'elle était inutile & \en{You needn't have cooked --- we ate out.} \textit{(Tu n'avais pas besoin de cuisiner.)} \tabularnewline
\end{tabularx}
\end{tableau}

\begin{astuce}[Les trois structures qui vous rendront immédiatement « naturel »]
\en{You should have told me!} (Tu aurais dû me le dire !) \quad
\en{It must have been hard.} (Ça a dû être dur.) \quad
\en{I could have done it.} (J'aurais pu le faire.)\\
Ces trois blocs, dits avec la bonne intonation, remplacent une heure de conversation
maladroite --- et expriment l'empathie, ce qui est le c\oe ur du contact humain.
\end{astuce}

\begin{exo}[1.1 --- Choisissez le modal]
\begin{enumerate}
  \item You \trou[1.8cm] smoke in the hospital. (interdiction)
  \item You \trou[1.8cm] pay, it's free. (absence d'obligation)
  \item \trou[1.8cm] I open the window? (permission polie)
  \item He \trou[1.8cm] be at home, the lights are on. (déduction forte)
  \item She \trou[1.8cm] come, but I'm not sure. (possibilité faible)
  \item You \trou[1.8cm] see a doctor. (conseil)
  \item When I was five, I \trou[1.8cm] already swim. (capacité passée)
  \item I \trou[1.8cm] to work every day --- it's the law. (obligation extérieure)
  \item I \trou[2.5cm] have called you yesterday, I'm sorry. (regret)
  \item They \trou[2.5cm] have missed the train --- they aren't here. (déduction passée)
\end{enumerate}
\end{exo}
\begin{corrige}[1.1]
1. \en{mustn't} 2. \en{don't have to} 3. \en{May / Can / Could} 4. \en{must}
5. \en{might / may / could} 6. \en{should} 7. \en{could} 8. \en{have}
9. \en{should have} 10. \en{must have}
\end{corrige}

\section{Unité 2 --- Les conditionnels}

\begin{tableau}[Les quatre conditionnels en un tableau]
\footnotesize
\begin{longtable}{@{}p{2.4cm}p{4.6cm}p{4.2cm}p{2.8cm}@{}}
\entete{Type} & \entete{Structure} & \entete{Exemple} & \entete{Emploi} \tabularnewline
\endfirsthead
\entete{Type} & \entete{Structure} & \entete{Exemple} & \entete{Emploi} \tabularnewline
\endhead
\hline
\textbf{0} & \en{if} + présent, présent & \en{If you heat water, it boils.} & vérité générale, fait scientifique \tabularnewline
\textbf{1} & \en{if} + présent, \en{will} + base & \en{If it rains, I'll stay home.} & situation réelle et probable dans l'avenir \tabularnewline
\textbf{2} & \en{if} + passé simple, \en{would} + base & \en{If I had time, I would travel.} & hypothèse improbable ou imaginaire au présent \tabularnewline
\textbf{3} & \en{if} + \en{past perfect}, \en{would have} + participe passé & \en{If I had studied, I would have passed.} & regret sur le passé, irréel \tabularnewline
\textbf{Mixte} & \en{if} + \en{past perfect}, \en{would} + base & \en{If I had taken that job, I would be rich now.} & cause passée, conséquence présente \tabularnewline
\end{longtable}
\end{tableau}

\begin{regle}[Le point qui fait toute la différence]
Dans la proposition en \en{if}, on emploie un temps \textbf{d'un cran en arrière} par rapport
au sens réel : présent pour du futur (type 1), passé pour du présent (type 2),
\en{past perfect} pour du passé (type 3). \textbf{Jamais de \en{will} après \en{if}}
(sauf dans \en{if you will} au sens de « si vous voulez bien »).
\end{regle}

\begin{piege}[\en{*If I will have time}]
\ko \en{If I will have time, I will call you.}\\
\ok \en{If I \textbf{have} time, I will call you.}\\
Même règle après \en{when, as soon as, unless, until, before, after} :
\en{When she arrives, we'll eat.} (pas \en{*when she will arrive}).
\end{piege}

\begin{astuce}[\en{If I were you\ldots} --- la formule de conseil]
\en{If I were you, I would take the job.} (Si j'étais toi, je prendrais le poste.)\\
C'est LE bloc à connaître pour donner un conseil avec élégance. \en{were} est la forme
correcte au subjonctif (toutes les personnes) ; \en{was} est accepté à l'oral.
\end{astuce}

\subsection{Les alternatives à \en{if}}
\begin{tableau}[Autres conjonctions conditionnelles]
\begin{tabularx}{\linewidth}{@{}p{3.4cm}p{4.4cm}X@{}}
\entete{Conjonction} & \entete{Sens} & \entete{Exemple} \tabularnewline
\hline
\en{unless} & à moins que (= \en{if\ldots not}) & \en{Unless you hurry, you'll be late.} \textit{(À moins que tu te dépêches, tu seras en retard.)} \tabularnewline
\en{as long as / so long as} & à condition que & \en{You can go, as long as you're back by ten.} \tabularnewline
\en{provided that / providing that} & à condition que (formel) & \en{We'll sign, provided that you agree.} \tabularnewline
\en{in case} & au cas où (précaution, \textbf{pas} une condition) & \en{Take an umbrella in case it rains.} \textit{(Prends un parapluie au cas où il pleuvrait.)} \tabularnewline
\en{otherwise / or} & sinon & \en{Hurry up, otherwise you'll miss it.} \tabularnewline
\en{even if} & même si & \en{Even if it rains, we'll go.} \tabularnewline
\en{suppose / supposing / what if} & et si\ldots ? & \en{Suppose you lost your job, what would you do?} \tabularnewline
\end{tabularx}
\end{tableau}

\begin{piege}[\en{in case} n'est pas « si »]
\en{I'll take an umbrella \textbf{in case} it rains.} = au cas où il pleuvrait (je me prépare).\\
\en{I'll take an umbrella \textbf{if} it rains.} = s'il pleut (condition).\\
Le français confond souvent les deux ; l'anglais les distingue nettement.
\end{piege}

\subsection{L'inversion conditionnelle (registre soutenu)}
\ex{Had I known, I would have come.}{Si j'avais su, je serais venu. (= \en{If I had known\ldots})}
\ex{Should you need help, please call me.}{Si vous avez besoin d'aide, appelez-moi. (= \en{If you should need\ldots})}
\ex{Were I richer, I would buy a house.}{Si j'étais plus riche, j'achèterais une maison.}

\begin{exo}[2.1 --- Mettez les verbes au bon temps]
\begin{enumerate}
  \item If it (rain) \trou[2cm] tomorrow, we (cancel) \trou[2.5cm] the match.
  \item If I (be) \trou[2cm] you, I (not / say) \trou[2.5cm] that.
  \item If she (study) \trou[2.5cm] harder, she (pass) \trou[2.5cm] the exam last year.
  \item I (buy) \trou[2cm] a car if I (have) \trou[2cm] enough money.
  \item If you heat ice, it (melt) \trou[2cm].
  \item Unless you (hurry) \trou[2cm], you (miss) \trou[2cm] the bus.
  \item If I (know) \trou[2.5cm] about the party, I (come) \trou[2.5cm].
  \item What (you / do) \trou[2.5cm] if you (win) \trou[2cm] the lottery?
  \item If he (not / forget) \trou[2.5cm] his passport, he (not / miss) \trou[2.5cm] the flight.
  \item Take a coat in case it (get) \trou[2cm] cold.
\end{enumerate}
\end{exo}
\begin{corrige}[2.1]
1. \en{rains} --- \en{will cancel} 2. \en{were / was} --- \en{wouldn't say}
3. \en{had studied} --- \en{would have passed} 4. \en{would buy} --- \en{had}
5. \en{melts} 6. \en{hurry} --- \en{will miss} 7. \en{had known} --- \en{would have come}
8. \en{would you do} --- \en{won} 9. \en{hadn't forgotten} --- \en{wouldn't have missed}
10. \en{gets}
\end{corrige}

\section{Unité 3 --- \en{wish}, \en{if only}, \en{would rather}, \en{it's time}}

\begin{tableau}[Exprimer le regret et le souhait]
\begin{tabularx}{\linewidth}{@{}p{3.6cm}p{4.2cm}X@{}}
\entete{Structure} & \entete{Sens} & \entete{Exemple} \tabularnewline
\hline
\en{wish} + passé simple & regret sur le \textbf{présent} & \en{I wish I \textbf{had} a car.} \textit{(J'aimerais avoir une voiture --- je n'en ai pas.)} \tabularnewline
\en{wish} + \en{past perfect} & regret sur le \textbf{passé} & \en{I wish I \textbf{had studied} more.} \textit{(J'aurais aimé étudier plus.)} \tabularnewline
\en{wish} + \en{would} & agacement, désir de changement chez l'autre & \en{I wish you \textbf{would stop} talking.} \textit{(J'aimerais que tu arrêtes de parler.)} \tabularnewline
\en{if only} + passé / \en{past perfect} & même sens, plus fort & \en{If only I \textbf{had} more time!} \tabularnewline
\en{would rather} + base & préférence (même sujet) & \en{I'd rather stay home.} \textit{(Je préférerais rester.)} \tabularnewline
\en{would rather} + sujet + passé & préférence (autre sujet) & \en{I'd rather you \textbf{didn't} smoke.} \textit{(Je préférerais que tu ne fumes pas.)} \tabularnewline
\en{it's (high) time} + passé & il serait temps que & \en{It's time we \textbf{left}.} \textit{(Il est temps qu'on parte.)} \tabularnewline
\en{hope} + présent / futur & espoir réel (pas de regret) & \en{I hope you \textbf{pass} the exam.} \textit{(J'espère que tu réussiras.)} \tabularnewline
\end{tabularx}
\end{tableau}

\begin{piege}[\en{wish} contre \en{hope}]
\en{I \textbf{hope} you feel better.} = j'espère que tu iras mieux (possible).\\
\en{I \textbf{wish} you felt better.} = j'aimerais que tu ailles mieux (mais ce n'est pas le cas).\\
\en{wish} implique toujours une dose d'irréel ou de regret. \en{hope}, lui, est tourné vers
l'avenir et s'emploie avec un temps normal. \en{*I wish you will come} est incorrect :
\en{I hope you will come} \ok\ ou \en{I wish you would come} \ok.
\end{piege}

\begin{exo}[3.1 --- Complétez]
\begin{enumerate}
  \item I wish I (be) \trou[2cm] taller. \item She wishes she (not / say) \trou[2.5cm] that.
  \item I wish it (stop) \trou[2cm] raining! \item I'd rather you (not / tell) \trou[2.5cm] anyone.
  \item It's time we (go) \trou[2cm] home. \item I hope you (have) \trou[2cm] a great trip.
  \item If only I (can) \trou[2.5cm] speak English perfectly.
  \item He wishes he (take) \trou[2.5cm] the other job last year.
\end{enumerate}
\end{exo}
\begin{corrige}[3.1]
1. \en{were / was} 2. \en{hadn't said} 3. \en{would stop} 4. \en{didn't tell}
5. \en{went} 6. \en{have / will have} 7. \en{could} 8. \en{had taken}
\end{corrige}

\section{Unité 4 --- La voix passive}

\begin{regle}[Formation]
\en{be} (au temps voulu) + \textbf{participe passé}. L'auteur de l'action passe au second plan
(avec \en{by}) ou disparaît.\\
\textbf{Actif :} \en{Somebody stole my bike.} \quad \textbf{Passif :} \en{My bike \textbf{was
stolen} (by somebody).}
\end{regle}

\begin{tableau}[Le passif à tous les temps]
\begin{tabularx}{\linewidth}{@{}p{3.6cm}p{5.0cm}X@{}}
\entete{Temps} & \entete{Passif} & \entete{Exemple} \tabularnewline
\hline
Présent simple & \en{am/is/are} + p.p. & \en{English is spoken here.} \tabularnewline
Présent continu & \en{am/is/are being} + p.p. & \en{The road is being repaired.} \tabularnewline
Passé simple & \en{was/were} + p.p. & \en{The house was built in 1920.} \tabularnewline
Passé continu & \en{was/were being} + p.p. & \en{The car was being washed.} \tabularnewline
\en{Present perfect} & \en{have/has been} + p.p. & \en{My phone has been stolen.} \tabularnewline
\en{Past perfect} & \en{had been} + p.p. & \en{The work had been finished.} \tabularnewline
Futur & \en{will be} + p.p. & \en{The results will be published tomorrow.} \tabularnewline
Modaux & \en{can/must/should be} + p.p. & \en{This must be signed. It can be done.} \tabularnewline
\en{going to} & \en{going to be} + p.p. & \en{A new bridge is going to be built.} \tabularnewline
Infinitif & \en{to be} + p.p. & \en{It needs to be cleaned.} \tabularnewline
\en{-ing} & \en{being} + p.p. & \en{I hate being interrupted.} \tabularnewline
\end{tabularx}
\end{tableau}

\subsection{Quand employer le passif ?}
\begin{enumerate}
  \item L'auteur est \textbf{inconnu} : \en{My wallet was stolen.}
  \item L'auteur est \textbf{évident ou sans importance} : \en{The rubbish is collected every Monday.}
  \item On veut rester \textbf{neutre ou poli} : \en{Mistakes were made.} (personne n'est accusé)
  \item Le résultat compte plus que l'acteur : \en{The bridge was opened in 2010.}
  \item Registre formel, administratif, scientifique, journalistique.
\end{enumerate}

\begin{piege}[Le français adore « on », l'anglais adore le passif]
\en{On m'a volé.} \ra \en{I was robbed.} (\en{*One stole me})\\
\en{On parle anglais ici.} \ra \en{English is spoken here.}\\
\en{On dit que c'est vrai.} \ra \en{It is said to be true.} / \en{They say it's true.}\\
\en{On a annulé la réunion.} \ra \en{The meeting was cancelled.}\\
Le pronom \en{one} existe en anglais mais il est littéraire et rare : ne l'utilisez pas
comme équivalent de « on ».
\end{piege}

\subsection{Verbes à deux compléments et passif impersonnel}
\ex{They gave him a prize. \ra\ He was given a prize. \ra\ A prize was given to him.}{On lui a donné un prix.}
\ex{They say that he is rich. \ra\ He is said to be rich. \ra\ It is said that he is rich.}{On dit qu'il est riche.}
\ex{People believe that the economy will improve. \ra\ The economy is believed to be improving.}{On pense que l'économie va s'améliorer.}

\subsection{\en{get} + participe passé : le passif familier}
\en{get married, get lost, get hurt, get fired, get arrested, get dressed, get broken,
get elected, get used to} --- très fréquent à l'oral, légèrement moins formel que \en{be}.
\ex{He got fired last week.}{Il s'est fait virer la semaine dernière.}

\subsection{\en{have something done} : faire faire quelque chose par un autre}
\begin{regle}
\en{have} + objet + \textbf{participe passé} = faire réaliser une action par quelqu'un d'autre.
\end{regle}
\ex{I had my hair cut.}{Je me suis fait couper les cheveux. (par un coiffeur)}
\ex{We're having the house painted next week.}{Nous faisons peindre la maison la semaine prochaine.}
\ex{I need to get my phone repaired.}{Je dois faire réparer mon téléphone. (\en{get} = variante familière)}

\begin{exo}[4.1 --- Mettez au passif]
\begin{enumerate}
  \item \en{They build these cars in Germany.}
  \item \en{Somebody has taken my pen.}
  \item \en{The storm destroyed the roof.}
  \item \en{People speak French in Senegal.}
  \item \en{They are repairing the road now.}
  \item \en{The company will launch a new product.}
  \item \en{You must sign this document.}
  \item \en{Someone told me the news.}
  \item \en{They say she is a genius.}
  \item \en{Somebody stole my bag on the bus.}
\end{enumerate}
\end{exo}
\begin{corrige}[4.1]
1. \en{These cars are built in Germany.} 2. \en{My pen has been taken.}
3. \en{The roof was destroyed (by the storm).} 4. \en{French is spoken in Senegal.}
5. \en{The road is being repaired now.} 6. \en{A new product will be launched.}
7. \en{This document must be signed.} 8. \en{I was told the news.}
9. \en{She is said to be a genius.} / \en{It is said that she is a genius.}
10. \en{My bag was stolen on the bus.}
\end{corrige}

\begin{exo}[4.2 --- \en{have something done}]
Traduisez :
\begin{enumerate}
  \item Je fais réparer ma voiture chaque année. \item Elle s'est fait voler son téléphone.
  \item Nous devons faire peindre la maison. \item Je me suis fait examiner les yeux hier.
\end{enumerate}
\end{exo}
\begin{corrige}[4.2]
1. \en{I have my car repaired every year.} 2. \en{She had her phone stolen.}
3. \en{We need to have the house painted.} 4. \en{I had my eyes tested yesterday.}
\end{corrige}

\section{Unité 5 --- Le discours rapporté (\en{reported speech})}

\begin{regle}[Le principe du « pas en arrière »]
Quand on rapporte des paroles au passé (\en{He said that\ldots}), le temps de la phrase
d'origine recule d'un cran, et les repères de temps et de lieu changent.
\end{regle}

\begin{tableau}[La transformation des temps]
\begin{tabularx}{\linewidth}{@{}p{5.4cm}X@{}}
\entete{Discours direct} & \entete{Discours rapporté} \tabularnewline
\hline
présent simple \ra & passé simple \tabularnewline
présent continu \ra & passé continu \tabularnewline
\en{present perfect} \ra & \en{past perfect} \tabularnewline
passé simple \ra & \en{past perfect} \tabularnewline
\en{will} \ra & \en{would} \tabularnewline
\en{can} \ra & \en{could} \tabularnewline
\en{may} \ra & \en{might} \tabularnewline
\en{must} / \en{have to} \ra & \en{had to} \tabularnewline
\en{should, could, might, would} \ra & inchangés \tabularnewline
\end{tabularx}
\end{tableau}

\begin{tableau}[Les changements de repères]
\begin{tabularx}{\linewidth}{@{}p{4.4cm}X@{}}
\entete{Direct} & \entete{Rapporté} \tabularnewline
\hline
\en{today / tonight} & \en{that day / that night} \tabularnewline
\en{yesterday} & \en{the day before / the previous day} \tabularnewline
\en{tomorrow} & \en{the next day / the following day} \tabularnewline
\en{now} & \en{then} \tabularnewline
\en{last week} & \en{the week before} \tabularnewline
\en{next month} & \en{the following month} \tabularnewline
\en{here} & \en{there} \tabularnewline
\en{this / these} & \en{that / those} \tabularnewline
\en{ago} & \en{before} \tabularnewline
\end{tabularx}
\end{tableau}

\subsection{\en{say} contre \en{tell}}
\begin{regle}
\en{say} + (proposition) : \en{He said (that) he was tired.}\\
\en{tell} + \textbf{personne} + (proposition) : \en{He told \textbf{me} (that) he was tired.}\\
Blocs figés : \en{tell the truth, tell a lie, tell a story, tell the time} ---
\en{say hello, say goodbye, say something, say a word}.
\end{regle}

\subsection{Rapporter une question}
\begin{regle}
\textbf{Question ouverte} (\en{wh-}) : on garde le mot interrogatif, puis ordre affirmatif.\\
\en{``Where do you live?''} \ra \en{He asked me where I lived.}\\
\textbf{Question fermée} : on introduit par \en{if} ou \en{whether}.\\
\en{``Are you married?''} \ra \en{She asked me if I was married.}
\end{regle}

\begin{piege}[\en{*He asked me where did I live}]
La faute la plus fréquente du discours rapporté : l'ordre interrogatif ne survit pas.
Après \en{asked}, c'est \textbf{sujet + verbe} : \en{where I lived}, \en{what she wanted},
\en{whether they had finished}.
\end{piege}

\subsection{Rapporter un ordre, une demande, un conseil}
\begin{regle}
Verbe rapporteur + \textbf{personne} + \en{to} + base (ou \en{not to} + base).
\end{regle}
\ex{``Close the door.'' \ra\ He told me to close the door.}{« Ferme la porte. » \ra Il m'a dit de fermer la porte.}
\ex{``Don't touch it.'' \ra\ She warned us not to touch it.}{« N'y touchez pas. » \ra Elle nous a prévenus de ne pas y toucher.}
\ex{``Please help me.'' \ra\ He asked me to help him.}{« Aide-moi, s'il te plaît. » \ra Il m'a demandé de l'aider.}

\subsection{Les verbes rapporteurs qui remplacent \en{say} (niveau B2)}
\begin{tableau}[Verbes rapporteurs et leurs structures]
\footnotesize
\begin{tabularx}{\linewidth}{@{}X X@{}}
\en{admit doing / that} --- admettre & \en{deny doing / that} --- nier \\
\en{promise to do} --- promettre & \en{refuse to do} --- refuser \\
\en{offer to do} --- proposer & \en{agree to do} --- accepter \\
\en{decide to do} --- décider & \en{plan to do} --- prévoir \\
\en{suggest doing / that} --- suggérer & \en{recommend doing} --- recommander \\
\en{apologise for doing} --- s'excuser de & \en{insist on doing} --- insister pour \\
\en{accuse someone of doing} --- accuser de & \en{blame someone for} --- reprocher à \\
\en{warn someone not to do} --- avertir de ne pas & \en{advise someone to do} --- conseiller de \\
\en{remind someone to do} --- rappeler de & \en{encourage someone to do} --- encourager à \\
\en{complain about} --- se plaindre de & \en{explain that / how} --- expliquer que \\
\end{tabularx}
\end{tableau}

\ex{``I'm sorry I'm late.'' \ra\ He apologised for being late.}{« Je suis désolé d'être en retard. » \ra Il s'est excusé d'être en retard.}
\ex{``You should see a doctor.'' \ra\ She advised me to see a doctor.}{« Tu devrais voir un médecin. » \ra Elle m'a conseillé de voir un médecin.}
\ex{``Don't forget to call me.'' \ra\ He reminded me to call him.}{« N'oublie pas de m'appeler. » \ra Il m'a rappelé de l'appeler.}
\ex{``Let's go out.'' \ra\ She suggested going out.}{« Sortons. » \ra Elle a proposé de sortir.}

\begin{astuce}[\en{that} est facultatif --- et souvent omis]
\en{He said (that) he was tired.} À l'oral, les natifs suppriment \en{that} la plupart du temps.
Le supprimer vous rend immédiatement plus naturel.
\end{astuce}

\begin{exo}[5.1 --- Rapportez ces phrases]
\begin{enumerate}
  \item \en{``I am very busy today.''} (he said)
  \item \en{``We have finished the work.''} (they told me)
  \item \en{``Will you help me?''} (she asked him)
  \item \en{``Where does she live?''} (I asked)
  \item \en{``Don't open the window.''} (the teacher told us)
  \item \en{``I will call you tomorrow.''} (he promised me)
  \item \en{``I saw her yesterday.''} (she said)
  \item \en{``Can you swim?''} (my mother asked me)
  \item \en{``You must leave now.''} (the guard told them)
  \item \en{``I'm sorry for the mistake.''} (he said)
\end{enumerate}
\end{exo}
\begin{corrige}[5.1]
1. \en{He said (that) he was very busy that day.}
2. \en{They told me (that) they had finished the work.}
3. \en{She asked him if / whether he would help her.}
4. \en{I asked where she lived.}
5. \en{The teacher told us not to open the window.}
6. \en{He promised to call me the next day.} / \en{He promised me that he would call me the following day.}
7. \en{She said (that) she had seen her the day before.}
8. \en{My mother asked me if / whether I could swim.}
9. \en{The guard told them (that) they had to leave then.}
10. \en{He apologised for the mistake.}
\end{corrige}

\begin{memo}[Chapitre 10 --- l'essentiel en 12 lignes]
\begin{enumerate}
  \item Après un modal : base verbale, jamais \en{to}, jamais \en{-s}.
  \item \en{mustn't} = interdit ; \en{don't have to} = pas obligé.
  \item \en{must be} = ça doit être ; \en{can't be} = c'est impossible.
  \item \en{should have done} = j'aurais dû.
  \item Type 1 : \en{if} + présent \ra \en{will}. Type 2 : \en{if} + passé \ra \en{would}.
  Type 3 : \en{if} + \en{past perfect} \ra \en{would have} + p.p.
  \item Jamais de \en{will} après \en{if}.
  \item \en{wish} + passé = regret présent ; \en{wish} + \en{past perfect} = regret passé.
  \item Passif = \en{be} + participe passé ; le français « on » se traduit souvent par le passif.
  \item \en{have something done} = faire faire par quelqu'un d'autre.
  \item Discours rapporté : un temps en arrière + changement des repères de temps.
  \item Après \en{asked}, ordre affirmatif : \en{where I lived} (pas \en{*where did I live}).
  \item Verbes rapporteurs riches : \en{promise to, admit doing, accuse of, advise to,
  apologise for, suggest doing}.
\end{enumerate}
\end{memo}
